\chapter{DF-027 --- Storybook-Style Tool for Bijou}\label{record-088}

\noindent{\footnotesize\color{BookMuted}\path{docs/design/DF-027-storybook-style-tool-for-bijou.md}\par}\smallskip

\begin{RecordMeta}
title: "DF-027 — Storybook-Style Tool for Bijou"
legend: DF
lane: design

\end{RecordMeta}

\section{Framing}\label{record-088-framing}

Bijou already has the pieces of a Storybook-style workflow: structured \texttt{C\allowbreak{}o\allowbreak{}m\allowbreak{}p\allowbreak{}o\allowbreak{}n\allowbreak{}e\allowbreak{}n\allowbreak{}t\allowbreak{}S\allowbreak{}t\allowbreak{}o\allowbreak{}r\allowbreak{}y} records, DOGFOOD's component explorer, story capture matrices, fixture-promotion records, and MCP-facing docs payloads. The missing piece is a standalone workstation that developers can use without entering the full docs app, backed by a stable model that tooling can consume without scraping DOGFOOD.

This cycle adds the first dedicated Storybook workstation slice. It includes an interactive TUI for local component browsing and development, plus the deterministic model/index and matrix capture path for agents, CI, MCP exports, block tooling, screenshots, and schema-bound controls.

\section{Sponsored Users}\label{record-088-sponsored-users}

\begin{itemize}
\tightlist
\item
  Bijou builders browsing first-party components and shells by family.
\item
  Maintainers auditing variant and output-mode coverage before release.
\item
  Tooling authors who need a stable story index and capture path for docs, MCP, screenshots, fixtures, or future block workstations.
\item
  Agents selecting stories without reverse-engineering the DOGFOOD shell.
\end{itemize}

\section{Hills}\label{record-088-hills}

\begin{enumerate}
\def\labelenumi{\arabic{enumi}.}
\tightlist
\item
  A builder can run \texttt{n\allowbreak{}p\allowbreak{}m\allowbreak{} \allowbreak{}r\allowbreak{}u\allowbreak{}n\allowbreak{} \allowbreak{}s\allowbreak{}t\allowbreak{}o\allowbreak{}r\allowbreak{}y\allowbreak{}b\allowbreak{}o\allowbreak{}o\allowbreak{}k} and enter a standalone Storybook-style workbench over the story catalog, with variants, profiles, package/source facts, and a preview surface.
\item
  A maintainer can capture a selected story across every variant/profile pair using the same \texttt{S\allowbreak{}t\allowbreak{}o\allowbreak{}r\allowbreak{}y\allowbreak{}C\allowbreak{}a\allowbreak{}p\allowbreak{}t\allowbreak{}u\allowbreak{}r\allowbreak{}e\allowbreak{}M\allowbreak{}a\allowbreak{}t\allowbreak{}r\allowbreak{}i\allowbreak{}x} substrate as existing docs tooling.
\item
  A future UI or MCP tool can consume a stable workstation model instead of scraping \texttt{e\allowbreak{}x\allowbreak{}a\allowbreak{}m\allowbreak{}p\allowbreak{}l\allowbreak{}e\allowbreak{}s\allowbreak{}/\allowbreak{}d\allowbreak{}o\allowbreak{}c\allowbreak{}s\allowbreak{}/\allowbreak{}a\allowbreak{}p\allowbreak{}p\allowbreak{}.\allowbreak{}t\allowbreak{}s} or re-parsing story docs.
\end{enumerate}

\section{Playback Questions}\label{record-088-playback-questions}

\begin{itemize}
\tightlist
\item
  Does the workstation index every DOGFOOD story without duplicating story truth?
\item
  Does it group stories by family and expose variants, profile modes, package, tags, and source path?
\item
  Can callers capture a selected story's variant/profile matrix?
\item
  Does the capture path render real story previews instead of placeholder text?
\item
  Does the tool reuse \texttt{c\allowbreak{}a\allowbreak{}p\allowbreak{}t\allowbreak{}u\allowbreak{}r\allowbreak{}e\allowbreak{}S\allowbreak{}t\allowbreak{}o\allowbreak{}r\allowbreak{}y\allowbreak{}M\allowbreak{}a\allowbreak{}t\allowbreak{}r\allowbreak{}i\allowbreak{}x\allowbreak{}(\allowbreak{})} rather than creating a second matrix format?
\item
  Is there a simple command for humans and agents to run?
\end{itemize}

\section{Requirements}\label{record-088-requirements}

\begin{itemize}
\tightlist
\item
  Add a standalone interactive Storybook workstation app.
\item
  Add a DOGFOOD storybook workstation module for deterministic tooling.
\item
  Add a text index renderer for the story catalog.
\item
  Add a selected-story matrix capture helper.
\item
  Reuse \texttt{C\allowbreak{}o\allowbreak{}m\allowbreak{}p\allowbreak{}o\allowbreak{}n\allowbreak{}e\allowbreak{}n\allowbreak{}t\allowbreak{}S\allowbreak{}t\allowbreak{}o\allowbreak{}r\allowbreak{}y} and \texttt{S\allowbreak{}t\allowbreak{}o\allowbreak{}r\allowbreak{}y\allowbreak{}C\allowbreak{}a\allowbreak{}p\allowbreak{}t\allowbreak{}u\allowbreak{}r\allowbreak{}e\allowbreak{}M\allowbreak{}a\allowbreak{}t\allowbreak{}r\allowbreak{}i\allowbreak{}x} contracts.
\item
  Add root scripts for the interactive workbench and text-first workstation.
\item
  Keep both paths on the same story catalog instead of duplicating story truth.
\end{itemize}

\section{Acceptance Criteria}\label{record-088-acceptance-criteria}

\begin{itemize}
\tightlist
\item
  \texttt{t\allowbreak{}e\allowbreak{}s\allowbreak{}t\allowbreak{}s\allowbreak{}/\allowbreak{}c\allowbreak{}y\allowbreak{}c\allowbreak{}l\allowbreak{}e\allowbreak{}s\allowbreak{}/\allowbreak{}D\allowbreak{}F\allowbreak{}-\allowbreak{}0\allowbreak{}2\allowbreak{}7\allowbreak{}/\allowbreak{}s\allowbreak{}t\allowbreak{}o\allowbreak{}r\allowbreak{}y\allowbreak{}b\allowbreak{}o\allowbreak{}o\allowbreak{}k\allowbreak{}-\allowbreak{}w\allowbreak{}o\allowbreak{}r\allowbreak{}k\allowbreak{}s\allowbreak{}t\allowbreak{}a\allowbreak{}t\allowbreak{}i\allowbreak{}o\allowbreak{}n\allowbreak{}.\allowbreak{}t\allowbreak{}e\allowbreak{}s\allowbreak{}t\allowbreak{}.\allowbreak{}t\allowbreak{}s} proves the cycle doc carries the modern playback sections.
\item
  The workstation model covers every \texttt{C\allowbreak{}O\allowbreak{}M\allowbreak{}P\allowbreak{}O\allowbreak{}N\allowbreak{}E\allowbreak{}N\allowbreak{}T\allowbreak{}\_\allowbreak{}S\allowbreak{}T\allowbreak{}O\allowbreak{}R\allowbreak{}I\allowbreak{}E\allowbreak{}S} entry.
\item
  The rendered index includes family grouping, counts, variants, modes, package names, and source paths.
\item
  The selected-story matrix captures every profile/variant combination and has no missing canonical modes for canonical DOGFOOD stories.
\item
  The matrix output contains real rendered story content.
\item
  \texttt{n\allowbreak{}p\allowbreak{}m\allowbreak{} \allowbreak{}r\allowbreak{}u\allowbreak{}n\allowbreak{} \allowbreak{}s\allowbreak{}t\allowbreak{}o\allowbreak{}r\allowbreak{}y\allowbreak{}b\allowbreak{}o\allowbreak{}o\allowbreak{}k} launches a standalone interactive Storybook workbench instead of deep-linking into \texttt{c\allowbreak{}r\allowbreak{}e\allowbreak{}a\allowbreak{}t\allowbreak{}e\allowbreak{}D\allowbreak{}o\allowbreak{}c\allowbreak{}s\allowbreak{}A\allowbreak{}p\allowbreak{}p\allowbreak{}(\allowbreak{})}.
\item
  \texttt{p\allowbreak{}a\allowbreak{}c\allowbreak{}k\allowbreak{}a\allowbreak{}g\allowbreak{}e\allowbreak{}.\allowbreak{}j\allowbreak{}s\allowbreak{}o\allowbreak{}n} exposes \texttt{n\allowbreak{}p\allowbreak{}m\allowbreak{} \allowbreak{}r\allowbreak{}u\allowbreak{}n\allowbreak{} \allowbreak{}s\allowbreak{}t\allowbreak{}o\allowbreak{}r\allowbreak{}y\allowbreak{}b\allowbreak{}o\allowbreak{}o\allowbreak{}k\allowbreak{}:\allowbreak{}i\allowbreak{}n\allowbreak{}d\allowbreak{}e\allowbreak{}x} for the deterministic text-first workstation, with \texttt{n\allowbreak{}p\allowbreak{}m\allowbreak{} \allowbreak{}r\allowbreak{}u\allowbreak{}n\allowbreak{} \allowbreak{}d\allowbreak{}o\allowbreak{}g\allowbreak{}f\allowbreak{}o\allowbreak{}o\allowbreak{}d\allowbreak{}:\allowbreak{}s\allowbreak{}t\allowbreak{}o\allowbreak{}r\allowbreak{}y\allowbreak{}b\allowbreak{}o\allowbreak{}o\allowbreak{}k} kept as a DOGFOOD-scoped compatibility alias.
\end{itemize}

\section{Implementation Outline}\label{record-088-implementation-outline}

\begin{enumerate}
\def\labelenumi{\arabic{enumi}.}
\tightlist
\item
  Add \texttt{e\allowbreak{}x\allowbreak{}a\allowbreak{}m\allowbreak{}p\allowbreak{}l\allowbreak{}e\allowbreak{}s\allowbreak{}/\allowbreak{}d\allowbreak{}o\allowbreak{}c\allowbreak{}s\allowbreak{}/\allowbreak{}s\allowbreak{}t\allowbreak{}o\allowbreak{}r\allowbreak{}y\allowbreak{}b\allowbreak{}o\allowbreak{}o\allowbreak{}k\allowbreak{}-\allowbreak{}a\allowbreak{}p\allowbreak{}p\allowbreak{}.\allowbreak{}t\allowbreak{}s}.
\item
  Add \texttt{e\allowbreak{}x\allowbreak{}a\allowbreak{}m\allowbreak{}p\allowbreak{}l\allowbreak{}e\allowbreak{}s\allowbreak{}/\allowbreak{}d\allowbreak{}o\allowbreak{}c\allowbreak{}s\allowbreak{}/\allowbreak{}s\allowbreak{}t\allowbreak{}o\allowbreak{}r\allowbreak{}y\allowbreak{}b\allowbreak{}o\allowbreak{}o\allowbreak{}k\allowbreak{}-\allowbreak{}w\allowbreak{}o\allowbreak{}r\allowbreak{}k\allowbreak{}s\allowbreak{}t\allowbreak{}a\allowbreak{}t\allowbreak{}i\allowbreak{}o\allowbreak{}n\allowbreak{}.\allowbreak{}t\allowbreak{}s}.
\item
  Export pure model/index/matrix helpers from that module.
\item
  Add a small CLI path that prints either the index or a selected story matrix.
\item
  Add a DF-027 cycle test around the model, text output, matrix capture, and script registration.
\item
  Move the backlog note into \texttt{d\allowbreak{}o\allowbreak{}c\allowbreak{}s\allowbreak{}/\allowbreak{}d\allowbreak{}e\allowbreak{}s\allowbreak{}i\allowbreak{}g\allowbreak{}n\allowbreak{}/} and update the v6 lane pointer.
\end{enumerate}

\section{Drift Check}\label{record-088-drift-check}

The current DOGFOOD app remains the canonical human-facing docs surface. This slice does not fork story data or add a browser dependency. The standalone workbench and text-first workstation both read the existing \texttt{C\allowbreak{}O\allowbreak{}M\allowbreak{}P\allowbreak{}O\allowbreak{}N\allowbreak{}E\allowbreak{}N\allowbreak{}T\allowbreak{}\_\allowbreak{}S\allowbreak{}T\allowbreak{}O\allowbreak{}R\allowbreak{}I\allowbreak{}E\allowbreak{}S} catalog and use the existing story protocol plus \texttt{c\allowbreak{}a\allowbreak{}p\allowbreak{}t\allowbreak{}u\allowbreak{}r\allowbreak{}e\allowbreak{}S\allowbreak{}t\allowbreak{}o\allowbreak{}r\allowbreak{}y\allowbreak{}M\allowbreak{}a\allowbreak{}t\allowbreak{}r\allowbreak{}i\allowbreak{}x\allowbreak{}(\allowbreak{})}.

The result is a narrow, reusable tooling seam: deterministic enough for agents and CI, but still powered by the same story records humans browse in DOGFOOD.

\section{Playback}\label{record-088-playback}

\begin{itemize}
\tightlist
\item
  RED: DF-027 existed only as a backlog concept, while DOGFOOD, showcase, MCP, and story capture each held partial storybook-shaped behavior.
\item
  GREEN: \texttt{c\allowbreak{}r\allowbreak{}e\allowbreak{}a\allowbreak{}t\allowbreak{}e\allowbreak{}D\allowbreak{}o\allowbreak{}g\allowbreak{}f\allowbreak{}o\allowbreak{}o\allowbreak{}d\allowbreak{}S\allowbreak{}t\allowbreak{}o\allowbreak{}r\allowbreak{}y\allowbreak{}b\allowbreak{}o\allowbreak{}o\allowbreak{}k\allowbreak{}W\allowbreak{}o\allowbreak{}r\allowbreak{}k\allowbreak{}b\allowbreak{}e\allowbreak{}n\allowbreak{}c\allowbreak{}h\allowbreak{}M\allowbreak{}o\allowbreak{}d\allowbreak{}e\allowbreak{}l\allowbreak{}(\allowbreak{})} produces a grouped story catalog from \texttt{C\allowbreak{}O\allowbreak{}M\allowbreak{}P\allowbreak{}O\allowbreak{}N\allowbreak{}E\allowbreak{}N\allowbreak{}T\allowbreak{}\_\allowbreak{}S\allowbreak{}T\allowbreak{}O\allowbreak{}R\allowbreak{}I\allowbreak{}E\allowbreak{}S}.
\item
  \texttt{r\allowbreak{}e\allowbreak{}n\allowbreak{}d\allowbreak{}e\allowbreak{}r\allowbreak{}D\allowbreak{}o\allowbreak{}g\allowbreak{}f\allowbreak{}o\allowbreak{}o\allowbreak{}d\allowbreak{}S\allowbreak{}t\allowbreak{}o\allowbreak{}r\allowbreak{}y\allowbreak{}b\allowbreak{}o\allowbreak{}o\allowbreak{}k\allowbreak{}I\allowbreak{}n\allowbreak{}d\allowbreak{}e\allowbreak{}x\allowbreak{}(\allowbreak{})} prints counts, required modes, families, stories, variants, packages, and source paths.
\item
  \texttt{c\allowbreak{}a\allowbreak{}p\allowbreak{}t\allowbreak{}u\allowbreak{}r\allowbreak{}e\allowbreak{}D\allowbreak{}o\allowbreak{}g\allowbreak{}f\allowbreak{}o\allowbreak{}o\allowbreak{}d\allowbreak{}S\allowbreak{}t\allowbreak{}o\allowbreak{}r\allowbreak{}y\allowbreak{}b\allowbreak{}o\allowbreak{}o\allowbreak{}k\allowbreak{}M\allowbreak{}a\allowbreak{}t\allowbreak{}r\allowbreak{}i\allowbreak{}x\allowbreak{}(\allowbreak{})} renders a selected story through every profile/variant pair and returns a \texttt{S\allowbreak{}t\allowbreak{}o\allowbreak{}r\allowbreak{}y\allowbreak{}C\allowbreak{}a\allowbreak{}p\allowbreak{}t\allowbreak{}u\allowbreak{}r\allowbreak{}e\allowbreak{}M\allowbreak{}a\allowbreak{}t\allowbreak{}r\allowbreak{}i\allowbreak{}x}.
\item
  \texttt{c\allowbreak{}r\allowbreak{}e\allowbreak{}a\allowbreak{}t\allowbreak{}e\allowbreak{}S\allowbreak{}t\allowbreak{}o\allowbreak{}r\allowbreak{}y\allowbreak{}b\allowbreak{}o\allowbreak{}o\allowbreak{}k\allowbreak{}A\allowbreak{}p\allowbreak{}p\allowbreak{}(\allowbreak{})} launches a standalone interactive TUI over the same story catalog, with catalog navigation, profile/variant switching, preview, and test-matrix panes.
\item
  \texttt{n\allowbreak{}p\allowbreak{}m\allowbreak{} \allowbreak{}r\allowbreak{}u\allowbreak{}n\allowbreak{} \allowbreak{}s\allowbreak{}t\allowbreak{}o\allowbreak{}r\allowbreak{}y\allowbreak{}b\allowbreak{}o\allowbreak{}o\allowbreak{}k} now targets the standalone interactive workbench instead of opening DOGFOOD on the components page.
\item
  \texttt{n\allowbreak{}p\allowbreak{}m\allowbreak{} \allowbreak{}r\allowbreak{}u\allowbreak{}n\allowbreak{} \allowbreak{}s\allowbreak{}t\allowbreak{}o\allowbreak{}r\allowbreak{}y\allowbreak{}b\allowbreak{}o\allowbreak{}o\allowbreak{}k\allowbreak{}:\allowbreak{}i\allowbreak{}n\allowbreak{}d\allowbreak{}e\allowbreak{}x} provides the text-first command path, while \texttt{n\allowbreak{}p\allowbreak{}m\allowbreak{} \allowbreak{}r\allowbreak{}u\allowbreak{}n\allowbreak{} \allowbreak{}d\allowbreak{}o\allowbreak{}g\allowbreak{}f\allowbreak{}o\allowbreak{}o\allowbreak{}d\allowbreak{}:\allowbreak{}s\allowbreak{}t\allowbreak{}o\allowbreak{}r\allowbreak{}y\allowbreak{}b\allowbreak{}o\allowbreak{}o\allowbreak{}k} remains available for existing docs and users.
\end{itemize}

\section{Retrospective}\label{record-088-retrospective}

The right Storybook-style tool has two layers: a standalone interactive workbench for humans and a stable model/capture path for CI, MCP, release audits, agents, and future block tooling. Both layers stay useful because they share the same story records instead of minting a second registry.

